Este es el recorrido. Son once partes y unas dos horas, sin contar las preguntas.

La primera parte es la más larga porque explica cómo se hicieron las pruebas: de dónde salieron los programas, por qué primero con Docker y luego sin Docker, y qué se corrió en cada máquina. Sin eso, los resultados no se pueden juzgar.

Después vienen los resultados, en el orden del paper: corrección, rendimiento y latencia. En medio van los hallazgos propios de la réplica.

El reparto entre los tres es una sugerencia; se puede mover.
